\chapter{Guided Walkthroughs}

\section{Purpose}

This chapter gives concrete source-reading paths. Each walkthrough starts from
a visible API or example and follows the mechanism down through the runtime.
The goal is to make the code base approachable without flattening the design
into prose only.

\begin{notabene}[How to use these walkthroughs]
Keep the source tree open while reading. The paths are intentionally explicit.
Read the public function first, then follow the listed modules in order.
\end{notabene}

\section{Walkthrough: spawning one task}

Start with the smallest async unit: a task spawned onto one executor.

\begin{lstlisting}[language={},caption={Task spawn reading path}]
examples/executor_sleep.rs
src/executor.rs
src/executor/spawner.rs
src/executor/task.rs
src/executor/task_state.rs
src/executor/scheduler.rs
src/executor/task_set.rs
\end{lstlisting}

The public surface is \texttt{executor\_and\_spawner}. It returns one executor
and one spawner. The spawner creates tasks; the executor polls them.

The useful questions to ask while reading are:

\begin{enumerate}
    \item Where is the future boxed and pinned?
    \item Where does the task get its id and optional name?
    \item What prevents repeated wakes from putting the same task into the
          ready queue many times?
    \item What state is recorded before and after each poll?
    \item What happens if the executor is dropped before the task completes?
\end{enumerate}

The mechanism is deliberately local. A task spawned by a \texttt{Spawner} is
owned by that spawner's executor. Its waker requeues it there. Nothing in this
path migrates the task to another thread.

\section{Walkthrough: sleeping and timers}

The timer path starts with examples such as \filepath{executor_sleep.rs} and
\filepath{executor_timeout.rs}. Then read:

\begin{lstlisting}[language={},caption={Timer reading path}]
src/executor/future.rs
src/executor/timer.rs
src/executor/scheduler.rs
src/executor/driver.rs
\end{lstlisting}

\texttt{sleep} creates a future that registers a deadline. When polled before
the deadline, the future stores the task waker and returns
\texttt{Poll::Pending}. The scheduler keeps timer state separate from the ready
queues. On each loop iteration the executor checks expired timers and wakes the
registered tasks.

\texttt{timeout} is useful for learning because it demonstrates cancellation by
drop. If the deadline wins, the inner future is dropped. Any timer or readiness
registration owned by that future must clean itself up. This is why the runtime
has tests for drop behavior rather than only success paths.

\section{Walkthrough: readiness-driven TCP}

Start from the server helpers rather than the raw reactor:

\begin{lstlisting}[language={},caption={TCP readiness reading path}]
examples/async_tcp_server.rs
examples/async_tcp_scoped_server.rs
src/executor/tcp.rs
src/executor/unix_io.rs
src/executor/io_interest.rs
src/os/reactor.rs
\end{lstlisting}

The TCP helper shows the application shape: accept a connection, spawn a
handler, wait for handlers, and propagate errors. The lower layers show the
runtime shape: try a non-blocking operation, register interest on
\texttt{WouldBlock}, sleep in the OS reactor, and wake the task when readiness
is reported.

When reading \filepath{src/executor/tcp.rs}, notice that the accept loop and
the handler tasks are distinct. Grouped TCP helpers use this distinction:
the accept loop stays in the caller's current task, while accepted handler
tasks are spawned into an explicit scheduling group.

\section{Walkthrough: scheduling groups}

The fastest way to see scheduling groups is to run:

\begin{lstlisting}[language=bash,caption={Scheduling group examples}]
cargo run --example scheduling_group_demo -- 1
cargo run --example async_tcp_scheduling_groups
cargo run --example sharded_scheduling_groups
\end{lstlisting}

Then read:

\begin{lstlisting}[language={},caption={Scheduling group reading path}]
src/executor/scheduling_group.rs
src/executor/spawner.rs
src/executor/task_set.rs
src/executor/snapshot.rs
src/sharded_executor.rs
\end{lstlisting}

\texttt{scheduling\_group.rs} defines the public handle and its executor
ownership. \texttt{spawner.rs} validates that a group belongs to the executor
before creating a task in that group. \texttt{task\_set.rs} owns the queues and
weighted virtual runtime accounting. \texttt{snapshot.rs} resolves group names
into task snapshots for observability.

On the sharded runtime, the interesting extra step is runtime identity. A
\texttt{ShardedSchedulingGroup} is a bundle of executor-local groups, one per
shard. It must not be accepted by a different runtime even if names and ids
look similar.

\section{Walkthrough: cross-shard submission}

Start with:

\begin{lstlisting}[language={},caption={Cross-shard reading path}]
examples/sharded_submit.rs
examples/sharded_map_reduce.rs
src/sharded_executor.rs
src/sharded_executor/join.rs
src/executor/spawner.rs
\end{lstlisting}

A cross-shard submit is not a function call into borrowed remote state. It is
owned work sent to another executor shard. The submitted future is polled on
the target shard. Its owned output completes a join handle. The original task
wakes and resumes on its original shard.

While reading the submitter APIs, look for the \texttt{Send + 'static}
boundaries. They are not incidental. They mark where work can cross threads.

\section{Walkthrough: shard-local state}

Read \filepath{examples/shard_local.rs} first, then:

\begin{lstlisting}[language={},caption={Shard-local reading path}]
src/shard_local.rs
src/sharded_executor.rs
src/runtime.rs
\end{lstlisting}

\texttt{ShardLocal<T>} is the clearest example of the shared-nothing idea in
the async runtime. Each shard owns one \texttt{T}. Access is routed to the
owning shard. The closure receives \texttt{\&mut T} synchronously and returns
an owned result.

The important rule is negative: a reference to \texttt{T} must not become part
of a future that later awaits. The API design enforces this by giving local
access through synchronous closures rather than async callbacks.

\section{Walkthrough: \texttt{io\_uring} lifecycle}

This path is Linux-specific. When available, start with:

\begin{lstlisting}[language=bash,caption={io\_uring examples}]
SITAS_DOCKER_IO_URING=1 tools/linux-docker.sh cargo run --example os_uring
SITAS_DOCKER_IO_URING=1 tools/linux-docker.sh cargo run --example os_uring_abandon
\end{lstlisting}

Then read:

\begin{lstlisting}[language={},caption={io\_uring reading path}]
src/os/uring.rs
src/executor/uring.rs
src/executor/driver.rs
src/executor/tests.rs
\end{lstlisting}

The key distinction is ownership. A readiness future borrows a descriptor and
retries ordinary system calls. An \texttt{io\_uring} future submits owned
buffers to the kernel and must keep them alive until completion or safe
abandonment cleanup.

\filepath{os_uring_abandon.rs} is especially useful because success paths can
hide lifetime complexity. Abandonment forces the dispatcher to keep enough
state to clean up operations whose futures no longer exist.

\section{Walkthrough: reading observability output}

Run:

\begin{lstlisting}[language=bash,caption={Observability examples}]
cargo run --example sharded_observability
cargo run --example shard_local_worker_observability
\end{lstlisting}

The columns are a compact runtime story:

\begin{description}
    \item[\texttt{ready}] how many tasks are queued to be polled;
    \item[\texttt{tasks}] how many tasks are still unfinished;
    \item[\texttt{timers}] how many timer registrations exist;
    \item[\texttt{polls}] how many task polls the executor has performed;
    \item[\texttt{done}] how many tasks completed or were cancelled;
    \item[\texttt{age\_ms}] how long the task has existed;
    \item[\texttt{state\_ms}] how long it has been in the current state;
    \item[\texttt{wait}] the last known wait interest.
\end{description}

The output is not tracing. There is no event stream. It is repeated sampling of
owned snapshots. This is less powerful than a full console, but it is simple,
portable, and keeps runtime internals inspectable.

\section{Debugging checklist}

When something behaves unexpectedly, ask these questions in order:

\begin{enumerate}
    \item Which shard owns the state being mutated?
    \item Which executor polls the future?
    \item Is the task queued, polling, waiting, completed, or cancelled?
    \item If waiting, is it waiting for a timer, fd readiness, completion, or
          an unknown waker?
    \item Did a handle from another executor or runtime cross a boundary?
    \item If using \texttt{io\_uring}, who owns the buffer until completion?
    \item If shutdown hangs, which spawner, submitter, join handle, or scoped
          child is still alive?
\end{enumerate}

These questions match the runtime's own invariants. They are usually more
productive than starting with performance guesses.
